\section{Emerging Fields and Open Problems}

Having characterized code as an agent harness through its interfaces, mechanisms, and orchestration patterns, we now examine how this paradigm materializes in concrete application domains and what open problems it exposes. Across coding assistants, GUI/OS agents, scientific discovery, personalization, and embodied agents, code serves not only as a model output, but also as the operational substrate for state representation, action execution, memory, feedback, and governance. These domains make the promise of code-centric agentic systems tangible, while revealing a common set of unresolved challenges around evaluation, verification, safety, coordination, multimodal grounding, and harness evolution.

\begin{figure}[t]
    \centering
    \includegraphics[width=1.0\linewidth]{figures/applications.pdf}
    \caption{Overview of code as an agent harness across five emerging domains, including coding assistants, GUI/OS agents, scientific discovery, personalization, and embodied agents.}
    \label{fig:applications}
\end{figure}

\subsection{Emerging Fields and Tangible Applications}
This subsection surveys five application domains where code-as-harness systems have become especially visible. Code assistants operate over repositories, tests, development tools, and collaborative workflows; GUI and OS agents manipulate rendered interfaces through executable actions and programmatic checkers; scientific agents organize hypotheses, experiments, analyses, and laboratory protocols as executable pipelines; personalization agents adapt recommendation policies through structured user feedback and editable preference states; and embodied agents ground high-level intent in executable skills subject to physical constraints. Together, these domains show how code connects model outputs to real-world systems, and how the design of the surrounding harness shapes reliability, controllability, and long-horizon autonomy.

\subsubsection{Code Assistants}
Code assistants provide one of the clearest application domains where
code-centric agentic systems become operational. Early systems mainly supported
localized completion or single-turn code generation. Recent assistants instead
operate across repository-level workflows, where editing, tool use, validation,
and pull-request interaction form a closed-loop agent process. This shift is
reflected in research systems such as SWE-agent~\citep{yang2024swe} and
OpenHands~\citep{wang2024openhands}, as well as production-oriented platforms
such as Claude Code~\citep{claudecode2025}, Codex~\citep{codex2025}, GitHub
Copilot coding agents~\citep{copilotagent2025}, and
DeepAgents~\citep{deepagents2025}. In these systems, the assistant is no longer
a standalone code generator. It is embedded in a development environment where
repository state, tools, validation routines, and collaboration workflows
provide the operational context for action and feedback.

\paragraph{Repository-centered Workspace}
Modern code assistants operate over repositories rather than isolated code
snippets. Source files, tests, build scripts, dependency metadata, issues,
branches, and pull requests form a persistent workspace that the agent can
inspect, modify, and validate over multiple steps. This makes repository-level
assistance less a matter of placing relevant files in the prompt, and more a
matter of constructing a task-specific working view over a large and evolving
codebase. Systems such as RepoCoder~\cite{zhang2023repocoder}, CodexGraph~\cite{liu2024codexgraphbridginglargelanguage}, and AutoCodeRover~\cite{zhang2024autocoderover} address this
problem through repository indexing, dependency-aware retrieval, graph-based
code representations, and agentic localization before editing. In this sense, the repository
becomes the operational substrate on which code assistants plan, act, and
receive feedback.

\paragraph{Executable Development Harnesses.}
Executable development harnesses are becoming the runtime and control plane of
code assistants. Rather than exposing the model to a flat list of tools, recent
systems wrap it in a managed development loop that controls repository access,
file edits, command execution, approval boundaries, context isolation, logging,
and validation. This trend is visible in production systems: Claude Code
packages local terminal/IDE/browser coding into a tool-mediated loop with
editing, command execution, permissions, hooks, memory, and subagents; Codex and
GitHub Copilot coding agents move similar loops into managed cloud or
GitHub-native workspaces with sandboxes, branches, approvals, and auditable
pull-request outputs; and DeepAgents exposes planning, filesystem-backed state,
context management, code execution, and subagent delegation as reusable harness
components~\citep{claudecode2025,codex2025,deepagents2025,copilotagent2025}.
Such loops are increasingly mediated by open protocols such as the Model Context Protocol~\citep{anthropic2024mcp,hou2025model}, 
which standardize how harnesses expose tools, context, and resources to the model and enable cross-system tool reuse. 
In parallel, recent research treats the harness itself as an object of optimization rather than a fixed wrapper: AutoHarness~\citep{lou2026autoharness} synthesizes
code harnesses from environment feedback, Meta-Harness~\citep{lee2026metaharness} searches over harness
code using prior candidates and execution traces, Agentic Harness Engineering~\citep{lin2026agentic}
evolves coding-agent harnesses through observability, and Natural-Language
Agent Harnesses~\citep{pan2026natural} externalize roles, contracts, adapters, and state conventions
into editable harness specifications. Together, these developments suggest that practical progress in code assistants
is increasingly shaped not only by improvements in the base model, but also by
the surrounding execution runtime, including its sandbox, permissions, context
plumbing, telemetry, and verification hooks.

\paragraph{Execution Feedback as Grounded Verification}
A distinguishing property of code assistants is the availability of machine-checkable feedback: compiler diagnostics, test outcomes, linter warnings, and runtime traces.
Agentless~\cite{xia2024agentless} shows that a fault-localization and patch-generation pipeline guided by test execution achieves competitive results on SWE-bench~\cite{jimenez2024swebench} without elaborate agentic control.
RepairAgent~\cite{bouzenia2025repairagent} and Live-SWE-agent~\cite{xia2025live} extend this loop into autonomous program repair driven by test results, while AlphaCodium~\cite{ridnik2024alphacodium} demonstrates that test-driven flow engineering substantially improves competitive programming performance over single-shot prompting.
Execution thus converts each candidate edit from a textual hypothesis into a verifiable transformation of the program world.

\paragraph{Memory and Context Management at Repository Scale}
Repositories routinely exceed any plausible context window, forcing code assistants to maintain explicit, structured memory.
Retrieval-augmented completion~\cite{zhang2023repocoder}, graph-based code indexing~\cite{liu2024codexgraphbridginglargelanguage}, documentation-oriented agents such as RepoAgent~\cite{luo-etal-2024-repoagent}, and recent context-retrieval benchmarks such as ContextBench~\cite{li2026contextbench} instantiate the memory abstractions of \S\ref{sec:memory} with a code-specific twist: stored items such as functions, tests, traces, and retrieved issue contexts are themselves executable or directly tied to executable states, and can be re-run, checked, or localized rather than merely re-read.
Recent memory systems further extend this view by storing reusable agent procedures or repository experience as procedural and experiential memory~\cite{gaurav2025codemem,wang2026memgovern}.
This narrows the gap between memory and environment found in conventional agent architectures, and makes abstraction management particularly acute, since the assistant must select the right scale of code and experience to surface for a given subtask.

\paragraph{Developer Intent and Project Conventions as Latent State}
Beyond explicit repository state, practical coding assistants must reason about
latent developer intent and project conventions. A useful patch should not only
pass visible tests, but also align with the repository's architecture, coding
style, and internal API reuse, properties that recent work describes as the
\emph{organicity} of generated code~\cite{li2026learning}. Agents that ignore
these constraints can produce technically correct patches that maintainers still
reject~\cite{li2026learning,thillen2026codetaste}, while benchmark analyses show
that some seemingly solved SWE-bench issues rely on solution leakage in the
issue text rather than genuine intent inference~\cite{aleithan2024swe}. Coding
assistance is therefore a partially observable program world problem: files,
tests, and tool outputs provide observable state, while design rationales,
implicit constraints, and team conventions must be inferred from issue threads,
prior commits, code reviews, and accumulated interaction history. This extends
the belief state divergence studied in SyncMind from shared multi agent state to
individual agent and user alignment~\cite{Guo2025SyncMind}. Modeling this latent
state is essential for moving from functional code generation toward trustworthy
developer collaboration.

\paragraph{From Inline Completion to Autonomous SWE Agents}
The evolution of code assistants can be viewed as an expansion of the development harness around the model.
Early systems such as Codex-based completion~\cite{chen2021evaluating} and commercial assistants such as Copilot~\cite{peng2023copilot} rely on a lightweight IDE harness, where local context is surfaced, an inline suggestion is generated, and the developer remains the primary executor, verifier, and state manager.
Productivity~\cite{peng2023copilot} and usability~\cite{vaithilingam2022expectation,mozannar2022reading} studies show that even this lightweight harness matters, since the value of a suggestion depends on its alignment with the developer's evolving program state and intent.
At the autonomous end, systems such as SWE-agent, OpenHands, AutoCodeRover, and Agentless operate within a repository-level harness, shifting from isolated code generation to stateful inspection, editing, execution, and revision.

\paragraph{From Patch Generation to Software Lifecycle Participation}
Code assistants are also moving from isolated patch generation toward broader
software lifecycle participation. SWE-bench framed repository-level assistance
as an issue-to-patch task~\citep{jimenez2023swe}, while newer benchmarks such as
SWE-Lancer~\citep{miserendino2025swe} and SWE-Bench Pro~\citep{deng2025swe} evaluate longer-horizon, economically meaningful
software deliverables that span multiple files and require professional
engineering effort. Related benchmarks
such as Terminal-Bench~\citep{merrill2026terminal} and AppWorld~\citep{trivedi2024appworld} further reflect the same shift toward
interactive environments where agents must operate through commands, tools, and
executable application states~\citep{xie2024osworld,yao2025taubench}. In deployment, this trend appears as agents that work
inside persistent engineering workflows rather than static repository snapshots,
including pull-request review, CI/CD feedback, and production issue resolution
~\citep{tang2024codeagent,Baqar_2025}. At production scale,
LingmaAgent reports that an autonomously deployed issue-resolution agent at
Alibaba Cloud resolves 16.9\% of in-house issues fully autonomously and 43.3\%
with manual intervention~\citep{ma2025alibaba,li2026advances}. This suggests
that code assistants are becoming workflow participants, not merely patch
generators.

\paragraph{Multi-Agent Code Assistance and Shared Repositories}
At the upper end of the spectrum, code assistance increasingly takes a multi-agent form, with planner, coder, tester, and reviewer roles operating over a shared repository.
ChatDev~\cite{Qian2023ChatDev}, MetaGPT~\cite{Hong2023MetaGPT},
CodeAgent~\cite{zhang2024codeagent}, and METAL~\cite{li2025metal} show how role specialization combined with a shared executable artifact enables coordination patterns that single agents struggle to sustain over long horizons.
The repository, together with its tests and execution traces, becomes both the medium of communication and the convergence target, directly instantiating the shared program world of \S\ref{sec:mas}.
Concurrent edits, however, can silently invalidate assumptions held by other agents, exposing the world-state synchronization challenges discussed in the same section.

\paragraph{The Harness as a Distillation Surface}
A defining 2026 development is that production harnesses are no longer only deployment infrastructure; they are becoming a dominant source of training data for the next generation of code-assistant models.
Cursor's Composer is trained with continuous online reinforcement learning on real Cursor usage traces, tightening the loop between deployed agent behavior and model updates~\citep{cursor2025composer,cursor2025rtrl}.
OpenAI's codex-1 (an o3 derivative)~\citep{codex2025}, GPT-5-Codex~\citep{openai2025gpt5codexcard}, and GPT-5.1-Codex-Max~\citep{openai2026codexmax} are explicitly trained on long-horizon, multi-turn coding interactions that mirror the Codex harness loop, while Anthropic's internal Claude Code dogfooding contributes a similar feedback channel documented in their teams-using-Claude-Code whitepaper~\citep{anthropic2025teams}.
At the same time, the harness itself is becoming an explicit optimization object: AutoHarness~\citep{lou2026autoharness} synthesizes harness code with a smaller LLM that filters illegal actions, Agentic Harness Engineering~\citep{lin2026agentic} closes an observability-driven evolution loop over harness components, Meta-Harness~\citep{lee2026metaharness} formalizes joint model--harness optimization, and Live-SWE-agent~\citep{xia2025live} edits its own scaffolding at runtime---together suggesting that the boundary between ``the agent'' and ``the harness around the agent'' is becoming a learnable surface in its own right.


\paragraph{Open Challenges for Code-Assistant Harnesses}
The maturation of production harnesses surfaces several coding-specific open problems that complement the cross-domain agenda discussed in the next subsection.
First, verification beyond unit tests remains largely unsolved: the oracle-adequacy crisis exposed by PatchDiff~\citep{wang2025solved} and SWE-Bench++~\citep{anonymous2025swebenchpp}, the security-correctness gap addressed by Aardvark~\citep{openai2025aardvark} and Codex Security~\citep{openai2026codexsecurity}, and the organicity gap between functional and accepted patches~\citep{li2026learning,thillen2026codetaste} all point to a verifier surface that current harnesses underspecify.
Second, failure attribution in long-horizon agent loops is still immature: empirical studies such as ``Why do multi-agent systems fail?''~\citep{cemri2025whymas}, the Who\&When attribution dataset~\citep{zhang2025whoandwhen}, AgenTracer~\citep{agentracer2025}, and AgentDebug~\citep{zhu2025llm} report best step-level attribution accuracies in the 14--53\% range, suggesting that production harnesses lack the structured traces needed for principled debugging.
Third, safety governance of autonomous code execution requires capability-based primitives that remain rare in practice: Aethelgard's learned capability governor~\citep{anonymous2026aethelgard}, fault-tolerant transactional sandboxing~\citep{anonymous2025faultsandbox}, and Microsoft's Agent Governance Toolkit~\citep{microsoft2026governance} represent early steps toward enforcing least privilege under concurrent agent action.
Fourth, harness self-evolution at production scale---demonstrated only in narrow settings by AutoHarness, AHE, and Live-SWE-agent---raises stability and rollback questions absent from non-self-modifying harnesses.
Fifth, multi-agent state synchronization on live repositories generalizes the SyncMind belief-state divergence problem~\citep{Guo2025SyncMind} to settings where humans, autonomous agents, and CI systems concurrently mutate shared program state.
Finally, trust calibration in pair programming user experience remains an under studied human factors problem, including decisions about when to interrupt, when to checkpoint, when to delegate, and when to defer, despite its centrality to whether harness driven autonomy can be safely scaled to enterprise workflows.

Code assistants are thus the clearest production instantiation of code-centric agentic systems and the most demanding testbed for the harness-engineering discipline now emerging across industry and academia.


\subsubsection{GUI/OS Agents as a Program World}

Graphical user interfaces and operating systems constitute, perhaps more than any other tangible application of foundation-model agents, a \textit{program world} in the most literal sense: every observation an agent receives is the rendered output of executable code (HTML, CSS, layout XML, accessibility APIs, framebuffers driven by window managers), and every action it takes is a call into another piece of code (a DOM event, an \texttt{adb} shell command, a keystroke captured by the OS event loop, a Playwright script). For this reason, GUI/OS agents have become the canonical testbed for the central thesis that code is the unifying substrate through which perception, action, environment dynamics, and memory can be represented, executed, and verified. Below we develop this view systematically.


\paragraph{GUI/OS as a Partially Observable Program World}

We model a GUI/OS environment as a Partially Observable Markov Decision Process $\langle \mathcal{S}, \mathcal{A}, \mathcal{O}, T, R\rangle$ in which the latent state \textit{s} $\in \mathcal{S}$ is the full program state of one or more processes (a browser's full DOM and JavaScript heap, an Android emulator's Activity stack and content providers, a Linux VM's filesystem and window tree). The agent never observes \textit{s} directly; it observes \textit{o} $\in \mathcal{O}$, which in modern systems takes one of four code-defined forms: (i) a serialized DOM or HTML subtree as in WebArena and Mind2Web \citep{zhou2024webarenarealisticwebenvironment,deng2023mind2webgeneralistagentweb}; (ii) an accessibility tree (AXTree) exposed by Android's UIAutomator or by macOS/Windows accessibility APIs as in AndroidWorld and WindowsAgentArena, for example, adopted by AgentOccam \citep{rawles2025androidworlddynamicbenchmarkingenvironment,bonatti2024windowsagentarenaevaluating,yang2024agentoccam}; (iii) a screenshot annotated with bounding-box or Set-of-Mark coordinates, the representation adopted by SeeAct, WebVoyager, OSWorld, and most recent native models \citep{zheng2024gpt4visiongeneralistwebagent,he2024webvoyagerbuildingendtoendweb,xie2024osworldbenchmarkingmultimodalagents,yang2023setofmarkpromptingunleashesextraordinary}; or (iv) hybrid representations that interleave pixels, accessibility metadata and HTML, as in WebArena's BrowserGym observation space and in CogAgent's dual-resolution encoder \citep{drouin2024workarenacapablewebagents,hong2024cogagentvisuallanguagemodel}. The action space $\mathcal{A}$ is likewise code: a tuple $\langle action\_type, target, value\rangle$ that compiles either to a DOM/accessibility call (\texttt{element.click()}, \texttt{setText(node\_id, ``...'')}) or to OS-level keyboard/mouse primitives (\texttt{pyautogui.click(x,y)}, \texttt{xdotool key}). Crucially, the transition function $T$ is not learned but \textit{executed}: the browser engine, the Android runtime, or the host OS deterministically produces the next observation. Agents are commonly framed as human-like computer users: they perceive the visual interface, reason over the user instruction, and execute actions through the same graphical channel available to humans. The agent's policy $\pi(a|h)$ is therefore best thought of as a \textit{program synthesizer} that, conditioned on a history \textit{h}, emits the next snippet of executable code; the environment is the interpreter.


\paragraph{Code as a Bridge Between User Interfaces and GUI Agents} 
Recent works treat code as an intermediate interface between high-level model reasoning and low-level UI execution \citep{xie2024osworldbenchmarkingmultimodalagents, wang2025guiagentsfoundationmodels,
xu2024androidlabtrainingsystematicbenchmarking}. This interface provides two main advantages: First, it abstracts away noisy visual details, and creates a natural boundary between the model's semantic planning and the system's executable control layer. Second, it fuses the perception, action, and evaluation in to a single code-as-harness pipeline.

On the action side, this is the GUI specialization of the broader CodeAct paradigm \citep{wang2024executablecodeactionselicit}: rather than emitting JSON tool calls, agents emit Python or JavaScript snippets that compose primitives such as \texttt{click(x, y)}, \texttt{type(text)}, \texttt{scroll(dx, dy)}, \texttt{key(``Enter'')}, and arbitrary library calls (e.g., \texttt{requests}, \texttt{subprocess}, \texttt{selenium}). Cradle makes this explicit by having an LMM output executable Python that drives keyboard and mouse for any application, including AAA games, achieving generalization across previously unseen software through skill curation and self-reflection rather than task-specific APIs \citep{tan2024cradleempoweringfoundationagents}. WebArena, BrowserGym, and TheAgentCompany similarly expose Playwright-style code actions whose execution is the ground truth of progress \citep{zhou2024webarenarealisticwebenvironment,drouin2024workarenacapablewebagents,xu2025theagentcompanybenchmarkingllmagents}.

On the perception side, recent native GUI models such as SeeClick, CogAgent, Ferret-UI, OS-Atlas, ShowUI, Aria-UI, UGround, UI-TARS, and GUI-Libra treat grounding as a \textit{function from pixels to executable coordinates}, training large vision-language models to emit $(x, y)$ or \texttt{bbox} tokens that can be directly piped into an action API \citep{cheng2024seeclickharnessingguigrounding,hong2024cogagentvisuallanguagemodel,you2024ferretuigroundedmobileui,wu2024osatlasfoundationactionmodel,lin2024showuivisionlanguageactionmodelgui,yang2025ariauivisualgroundinggui,gou2025navigatingdigitalworldhumans,qin2025uitarspioneeringautomatedgui,yang2026guilibratrainingnativegui}. By collapsing the planner→grounder→executor pipeline into a single VLA model whose output token stream is itself runnable code, these systems eliminate the brittle string-matching layer that historically separated language plans from grounded actions, as documented in SeeAct's analysis showing that grounding, rather than planning, is the dominant bottleneck on Mind2Web \citep{zheng2024gpt4visiongeneralistwebagent}.

On the evaluation side, code-defined environments enable \textit{executable feedback}: success is determined not by a learned reward model but by running an evaluator script over the post-action system state. WebArena's URL/string assertions, OSWorld's per-task Python checkers operating over OS file I/O and application state, AndroidWorld's \texttt{adb}-based state inspection, and Spider2-V's enterprise-tool checks all share the same pattern, an evaluator is itself a piece of code that interrogates the program world after the agent has finished \citep{zhou2024webarenarealisticwebenvironment,xie2024osworldbenchmarkingmultimodalagents,rawles2025androidworlddynamicbenchmarkingenvironment,cao2024spider2vfarmultimodalagents}. This closes the loop: code generates the environment, code is the agent's action, and code adjudicates the result.

\paragraph{Memory as Persistent Programmatic State}
For code-grounded GUI agents, memory is best understood as a \textit{persistent programmatic state layer}: structured artifacts that outlive the current UI state and can be retrieved, composed, or executed in later interactions. Recent works explore different line of memory: (i) \textit{Working memory of UI state} compresses the current observation to a task-relevant abstraction: Synapse's state-abstraction module filters HTML to a few task-relevant elements, allowing trajectory-as-exemplar prompting and an exemplar memory that retrieves prior trajectories by similarity \citep{zheng2024synapsetrajectoryasexemplarpromptingmemory}. (ii) \textit{Long-term cross-app/session memory} is implemented as structured documents and skill libraries: AppAgent compiles an exploration document per application that records the learned function of each UI element, which is then consulted on subsequent tasks \citep{zhang2023appagentmultimodalagentssmartphone}; Mobile-Agent-v2 introduces a dedicated planning agent whose memory tracks long-horizon progress across sub-tasks \citep{wang2024mobileagentv2mobiledeviceoperation}; Cradle maintains an explicit skill-curation module that promotes successful code snippets to a reusable library \citep{tan2024cradleempoweringfoundationagents}. Whereas these designs are tightly coupled to the host application's UI ontology, PlugMem proposes a \textit{task-agnostic} plugin memory module that distils raw interaction traces into a compact knowledge-centric memory graph of propositional and prescriptive knowledge, transferring unchanged from web agents to long-horizon dialogue and multi-hop retrieval \citep{yang2026plugmemtaskagnosticpluginmemory}. (iii) \textit{Self-evolving GUI agents} (already cited in this survey as UI-Voyager \citep{lin2026uivoyagerselfevolvingguiagent}) and AutoGLM extend this idea with online curriculum reinforcement learning that continuously grows a library of grounded behaviors, while OS-Genesis and UI-TARS use reflective trace collection on hundreds of virtual machines as a form of distilled memory \citep{liu2024autoglmautonomousfoundationagents,sun2025osgenesisautomatingguiagent,qin2025uitarspioneeringautomatedgui}. In all three regimes the memory is itself a code artifact, for example, a JSON document, a Python skill module, or a vector index of code-formatted trajectories, directly executable or directly composable into the agent's next action.

\paragraph{UI Simulators and Sandboxes as Executable Dynamics}

The simulator stack for GUI/OS agents is perhaps the clearest demonstration that environment dynamics in this domain \textit{is} code. Early benchmarks such as MiniWoB++ defined each task as a self-contained HTML/JavaScript page with a programmatic reward function \citep{liu2018reinforcementlearningwebinterfaces}; WebShop scaled this to 1.18M real Amazon products inside a self-hosted shopping site \citep{yao2023webshopscalablerealworldweb}. Mind2Web cached real-world traces for offline evaluation, while WebArena and VisualWebArena fork four full-stack open-source sites into Docker containers with deterministic resets and per-task functional checkers \citep{deng2023mind2webgeneralistagentweb,zhou2024webarenarealisticwebenvironment,koh2024visualwebarenaevaluatingmultimodalagents}. OSWorld pushes this further to 369 real Ubuntu/Windows/macOS tasks in disposable VMs whose initial state, golden actions, and Python evaluation scripts are all version-controlled artifacts \citep{xie2024osworldbenchmarkingmultimodalagents}; WindowsAgentArena specializes the same architecture for Windows 11 with Azure-parallel execution \citep{bonatti2024windowsagentarenaevaluating}; and Spider2-V extends OSWorld to professional data-engineering pipelines spanning BigQuery, dbt, and Airbyte \citep{cao2024spider2vfarmultimodalagents}. On mobile, AndroidWorld provides 116 programmatic tasks dynamically parameterized from natural-language templates with reward signals derived from device system state, while AndroidArena and AndroidLab supply complementary cross-app evaluations \citep{rawles2025androidworlddynamicbenchmarkingenvironment,xing2024understandingweaknesslargelanguage,xu2024androidlabtrainingsystematicbenchmarking}. BrowserGym and WorkArena unify many of these under a common Gym-style API and add 23,150 enterprise ServiceNow task instances \citep{drouin2024workarenacapablewebagents}, while AgentBench's OS and web tracks and the OpenHands-driven TheAgentCompany benchmark situate GUI control inside broader knowledge-work simulations \citep{liu2025agentbenchevaluatingllmsagents,xu2025theagentcompanybenchmarkingllmagents}. Most recently, Code2World makes the program-world stance explicit at the model level by training a vision-language coder that predicts the next GUI state as renderable HTML, turning the world model itself into an executable artifact and using rendered outcomes as reinforcement signals \citep{zheng2026code2worldguiworldmodel}. Together, these sandboxes embody the survey's claim that environment dynamics in agentic systems are increasingly authored as code: they are forkable, diffable, version-controlled, and reproducible in ways that no learned simulator can match.

\paragraph{From Simulation to Production: Executable Feedback Loops}

The same code-as-harness interface that makes simulators tractable has enabled an unusually rapid jump to production deployment, because the agent's input/output contract: screenshots in, code (or coordinate-typed function calls) out, is identical in both settings. Anthropic's Claude Computer Use exposes a public-beta API in which the model takes screenshots of a sandboxed desktop and emits keyboard/mouse actions as structured tool calls \citep{anthropic2024computeruse}. OpenAI's Operator and the underlying Computer-Using Agent (CUA) followed, combining GPT-4o's vision with reinforcement-learned reasoning over a unified click/scroll/type action space \citep{openai2025operator}. Google DeepMind's Project Mariner ships a Gemini-powered Chrome extension that observes the rendered DOM, plans, and executes browser actions on behalf of the user, and is being integrated into Search's AI Mode and the Gemini app \citep{deepmind2025mariner}. ByteDance's UI-TARS-1.5/2 and the associated UI-TARS-desktop product, Zhipu's AutoGLM (web browser plug-in and Android app), and Tencent's AppAgent lineage demonstrate that the same architecture transfers from the lab to consumer devices \citep{qin2025uitarspioneeringautomatedgui,liu2024autoglmautonomousfoundationagents,zhang2023appagentmultimodalagentssmartphone}. AutoWebGLM, the production sibling of CogAgent, exemplifies the route from arXiv preprint to deployed browser agent through an ``intermediate interface'' that decouples planning from grounding \citep{lai2024autowebglmlargelanguagemodelbased}. Earlier industrial efforts, like Adept's ACT-1/ACT-2 and Rabbit's Large Action Model, anticipated this trajectory but predated the executable-feedback infrastructure that has since made the loop reliable enough for deployment.

Looking forward, the literature converges on three frontiers, all expressed in code-as-harness terms. \textit{First}, native end-to-end agents that internalize perception, planning, grounding, and action into a single VLA model are displacing the modular planner+grounder pipeline. \textit{Second}, executable world models promise to give agents human-like foresight by predicting the next UI state as renderable code rather than as pixels or unstructured text. \textit{Third}, embodied, instruction-following GUI agents treat the entire device (e.g., terminal, browser, native apps, and peripherals) as a unified program world. The common thread is that code is the lingua franca: it defines observations, actions, evaluation, memory, and increasingly the world model itself.


\subsubsection{Autonomous Embodied Agents}

Embodied agent operates in the physical world or its simulation, perceiving the environment through structured outputs from vision and force sensors, and acting through motor commands subject to physical constraints such as reachability, collision, and dynamics.

\paragraph{Code as the Control Boundary that Connecting Agents and the World}
Unlike purely reasoning agents, embodied agents operate under physical constraints that may fail silently when violated: a robot may attempt to grasp an object outside its workspace without producing any explicit failure signal~\citep{liang2023codepolicieslanguagemodel}. This shifts the burden of correctness from runtime to action-generation time, where the agent's output must already be expressive enough to compose verified operation intents before reaching the actuator. 
Code naturally satisfies the requirements by serving both as the grounding interface and as the safety boundary. As a grounding interface, it translates high-level intent from LLMs into embodiment-respecting commands through primitive skill calls \citep{ahn2022can, ren2023robots, zhai2026skillvla, zhang2023bootstrap}, synthesized Python control policies \citep{liang2023code, mu2024robocodex, xie2025robotic, wang2025llm, ji2026genswarm}, and structured behavior-tree programs \citep{zhang2025codebt}. As a safety boundary, it constrains admissible actions at execution time \citep{guan2025normcode, szeider2025cp, miculicich2025veriguard}.

\paragraph{Layered Harness for Grounded and Verifiable Embodied Actions} 
Embodied agents require a layered harness that separates semantic reasoning from executable, physically grounded, and human-governed control~\citep{vemprala2024chatgpt}. Foundation models handle the semantic layer of embodied agency: interpreting goals, decomposing tasks, inferring affordances, selecting skills, proposing actions, and replanning under changing observations~\citep{huang2022inner, wang2023voyager}.  Code and classical robotics software define the admissibility boundary by exposing typed robot APIs, parameterizing primitive skills, calling geometric libraries, invoking motion planners, and supporting inspection, replay, versioning, and verification \citep{xie2025robotic, liang2023code, huang2023voxposer, macenski2020nav2}. Perception models and state estimators convert raw sensor streams into structured state that planners and controllers can use ~\citep{driess2023palme, deepmind2025geminirobotics}. Physical systems and low-level controllers then enforce embodiment-specific constraints such as kinematics, dynamics, collision avoidance, workspace limits, contact forces, timing, and stability. 

\paragraph{Reusable Skills as Embodied Memory} While code grounds a single action in physical feasibility, embodied agents operating over long horizons must also accumulate experience across tasks. In this regime, code takes on a second role: the same executable form that makes an action verifiable also makes it storable and reusable. 
Memory therefore naturally takes the form of a skill library, a collection of code artifacts that record past behavior and can be called as actions in future tasks.
This dual identity distinguishes embodied memory from other memory abstractions in \S\ref{sec:memory}: a skill is not merely something the agent reads, but something the agent re-executes. Voyager pioneered this paradigm with an growing skill library for open-ended tasks in Minecraft \citep{wang2023voyager}, and other work extends the same idea along several directions: tabletop manipulation \citep{tziafas2024lifelong}, human correction \citep{meng2025growing}, vision-grounded replanning \citep{kagaya2025vireskill}, and continual learning \citep{wang2026lifelong}. The principle has even crossed into the GUI domain \citep{lin2026uivoyagerselfevolvingguiagent}. Across these systems, the challenge has shifted from generating skills to governing the library: handling forgetting, abstraction, and grounding alignment.

\paragraph{Coordinated and Auditable Real-World Deployment} Moving from simulation to real-world deployment introduces challenges that go beyond a single agent: multiple robots must coordinate, behaviors must be auditable, and skills must transfer across embodiments. Code naturally extends to address all three. For coordination, it provides the substrate for multi-robot policy synthesis \citep{ji2026genswarm} and robot-agnostic cooperative architectures \citep{ashley2026racas}. For auditability, it supports governance mechanisms for industrial safety \citep{guan2025normcode, liu2026agents4plc} and verified closed-loop control \citep{santos2026alrm}. For cross-embodiment transfer, the same code-based skill abstraction enables combinatorial reuse on dual-arm systems \citep{zhai2026skillvla}. Open challenges remain in reducing the sim-to-real gap, scaling multi-agent coordination, and maintaining safety as environments evolve.




\subsubsection{Agents for Scientific Discovery as Program Worlds}

Scientific research is among the most natural testbeds for code as an agent harness: the scientific method is itself a closed loop of \textit{hypothesize → design → execute → observe → revise}, in which each transition is mediated by an artifact that is, increasingly, a program. Modern science can already be digital end-to-end, for example, hypotheses are encoded as differential equations or generative models, experimental protocols are written as XDL or Opentrons scripts, instruments are driven through Python APIs, and analyses live in Jupyter notebooks whose cells form a verifiable trace of reasoning. This makes scientific discovery an ideal domain to instantiate the three-fold role of code: code as the medium of \textit{reasoning} (e.g., symbolic derivations, formal proofs, hypothesis-as-program), code as the substrate of \textit{acting} (e.g., calls to wet-lab robots, simulators, statistical pipelines), and code as the executable \textit{environment} itself (e.g., molecular-dynamics engines, autonomous laboratories, virtual research teams). Recent systems, like \textit{AI Scientist v1/v2}, \citep{lu2024aiscientistfullyautomated,yamada2025aiscientistv2workshoplevelautomated} AI co-scientist \citep{gottweis2025aicoscientist}, Virtual Lab \citep{swanson2025virtual} and Biomni \citep{huang2025biomni}, make this code-as-harness framing concrete by elevating the entire research workflow to a single, executable program graph.

\paragraph{Scientific Discovery as a Partially Observable Program World}

We treat a research project as a partially observable program world $\langle \mathcal{S}, \mathcal{A}, T, \mathcal{O}, R\rangle$. The state $\mathcal{S}$ is a structured program memory containing the current best hypotheses, accumulated literature, code artifacts, intermediate datasets, and experimental observations. Actions $\mathcal{A}$ are typed code expressions: literature-search queries, calls to symbolic or numerical solvers, generation of new experimental scripts, modifications to a training pipeline, or robot-control commands. The transition function $T$ is realized by a Python interpreter, a Lean kernel, a quantum-chemistry package, a robotic synthesizer, or, in fully end-to-end systems such as the AI Scientist v2 \citep{yamada2025aiscientistv2workshoplevelautomated}, by a tree-search experiment manager that orchestrates all of these. Observations $\mathcal{O}$ correspond to execution outputs (numerical results, plots, error messages, peer-review scores), and the latent reward $R$ encodes desiderata such as novelty, reproducibility, and statistical significance. Crucially, the policy of a scientific agent is itself a program: ChemCrow \citep{bran2023chemcrowaugmentinglargelanguagemodels} composes 18 expert-designed chemistry tools through structured tool calls; Coscientist \citep{boiko2023autonomous} interleaves Python execution, web search, and robotic-API actions; and \textit{AlphaProof} \citep{hubert2025olympiad} expresses each ``reasoning step'' as a Lean tactic that the proof assistant verifies before transitioning the state. This view recasts traditionally informal categories (e.g., hypothesis, protocol, claim) as concrete program objects whose execution traces can be logged, replayed, and audited.

\paragraph{Unifying Ideation, Experimentation, Analysis, and Communication}

Traditional accounts of science separate ideation, experiment design, data analysis, and dissemination into distinct workflows with distinct tools. Code-centric agents collapse these into a single executable pipeline. ResearchAgent \citep{baek2025researchagentiterativeresearchidea} and SciAgents-style systems iteratively refine hypotheses by traversing entity graphs over the literature, with each candidate idea materialized as a structured object that can be passed to downstream planners. BioPlanner \citep{odonoghue2023bioplannerautomaticevaluationllms} formalizes wet-lab protocols as pseudocode whose admissible functions can be type-checked, retrieved, and composed, providing the same compositional substrate for biology that XDL provides for chemistry \citep{mehr2020universal}. Agent Laboratory \citep{schmidgall2025agentlaboratoryusingllm} and its preprint-sharing extension AgentRxiv \citep{schmidgall2025agentrxivcollaborativeautonomousresearch} explicitly factor research into three program-level phases: literature review, experimentation, report writing, orchestrated by specialized PhD, postdoc, and engineer agents that exchange Python files, LaTeX, and arXiv records. The AI Scientist \citep{lu2024aiscientistfullyautomated,yamada2025aiscientistv2workshoplevelautomated} goes further by representing an entire ML paper as a single executable trace: the system writes the experimental code with a coding assistant, executes it, reads the figures with a vision-language model, and emits a LaTeX manuscript that includes the very plots it generated. In all of these systems, what used to be a heterogeneous pipeline of natural-language artifacts becomes a homogeneous flow of typed code objects, enabling end-to-end optimization and automatic verification at every stage \citep{ren2026scientificintelligencesurveyllmbased,wang2024executablecodeactionselicit}.

\paragraph{Memory as Persistent Program State}

Long-horizon research depends on memory: prior experiments, failed attempts, citation graphs, and tacit lab know-how. Code-centric agents externalize this memory as persistent program state. At the \textit{working-memory} level, agents maintain executable scratchpads,  typically a Jupyter kernel or a CodeAct-style Python REPL \citep{jiang2025aideaidrivenexplorationspace},  whose live variables, dataframes, and figures form the immediate context for reasoning. El Agente Q \citep{Zou_2025} and Biomni \citep{huang2025biomni} exemplify hierarchical memory: short-lived tool outputs are cached in an episodic buffer, while structured artifacts (plasmid maps, optimized geometries, fitted models) are written to durable file stores that subsequent agent steps can re-load. At the \textit{long-term} level, PaperQA / PaperQA2 \citep{lala2023paperqaretrievalaugmentedgenerativeagent} and Google's AI co-scientist \citep{gottweis2025aicoscientist} treat the scientific literature itself as an indexed knowledge base, accessed through tool calls that retrieve passages, expand citations, and detect contradictions; this enables hypothesis evaluation against millions of prior results without inflating the prompt. AgentRxiv \citep{schmidgall2025agentrxivcollaborativeautonomousresearch} takes the idea one step further by giving autonomous research agents a shared preprint server: hypotheses, code, and findings produced by one run are uploaded as durable program artifacts that future runs can build on, instantiating cumulative scientific progress as a globally shared, version-controlled program state. Biomni's action-discovery agent \citep{huang2025biomni} mines tens of thousands of bioRxiv papers to populate a unified tool registry across 25 biomedical subfields, so that ``remembering how to clone a plasmid'' becomes the concrete act of importing a verified, code-level protocol from persistent storage.

\paragraph{Simulators as Executable Dynamics}

Scientific agents rely on simulators of physical and computational reality, and the code-as-harness view treats these uniformly as executable transition models. In computational chemistry, El Agente Q \citep{Zou_2025} wraps DFT engines, geometry optimizers, and thermochemistry tools as callable functions that the LLM invokes to roll out alternative reaction trajectories; on six university-level benchmarks it exceeds 87\% task success while emitting a transparent action-trace log of every simulation. ChemCrow \citep{bran2023chemcrowaugmentinglargelanguagemodels} similarly integrates RDKit, retrosynthesis engines, and reaction predictors so that an agent can ``execute'' a candidate synthesis virtually before committing to a wet-lab run. In structural and systems biology, the Virtual Lab \citep{swanson2025virtual} composes ESM, AlphaFold-Multimer, and Rosetta into a Python pipeline through which an LLM Principal-Investigator agent and its subordinate scientist agents jointly designed 92 SARS-CoV-2 nanobodies, two of which showed validated binding to JN.1 and KP.3 variants,  all in a few days of simulated meetings. For algorithmic and mathematical science, AlphaProof \citep{hubert2025olympiad} uses the Lean theorem prover as the executable environment, formally verifying every candidate proof step before reinforcing the language model, and AlphaEvolve \citep{novikov2025alphaevolvecodingagentscientific} orchestrates an evolutionary loop in which Gemini-generated code edits are executed and scored by automated evaluators, yielding new matrix-multiplication algorithms and mathematical constructions. In each case the simulator is the world: program states evolve only through verified executions, eliminating much of the hallucination that plagues purely textual scientific reasoning \citep{ren2026scientificintelligencesurveyllmbased}.

\paragraph{From Simulation to Production: Self-Driving Labs as Executable Feedback Loops}

The decisive test of a scientific agent is whether its closed loop crosses the boundary into physical reality. Self-driving laboratories (SDLs) are the production systems of this domain: they expose real instruments, like liquid handlers, XRD scanners, spectrometers, robotic arms, through code APIs, and accept agent-generated programs as their primary input. Berkeley's A-Lab \citep{szymanski2023autonomous} combines machine-learned synthesis recipes with autonomous robotics to synthesize 41 novel inorganic compounds from a target list of 58 in 17 days of continuous operation, while early thin-film SDLs \citep{macleod2020selfdrivinglaboratoryaccelerateddiscovery} established that Bayesian optimization loops can be wrapped as Python services and run unattended. Coscientist \citep{boiko2023autonomous} crossed this threshold for organic chemistry by autonomously planning, executing, and analyzing palladium-catalyzed Suzuki and Sonogashira couplings on the Emerald Cloud Lab and an in-house liquid-handling platform from a single English prompt. The Cronin group's Chemputer and its XDL chemical-description language \citep{mehr2020universal} formalize this contract: any synthesis published in the literature can be parsed into hardware-independent XDL code that compiles, like LLVM IR for chemistry, onto any compliant robotic platform. In biology, Biomni \citep{huang2025biomni} generates end-to-end molecular-cloning protocols that human reviewers rated comparable to a senior Stanford postdoc, while Google's AI co-scientist's drug-repurposing and antimicrobial-resistance hypotheses were experimentally validated in collaborator wet labs at Imperial College and Stanford \citep{gottweis2025aicoscientist}. MatPilot \citep{ni2024matpilotllmenabledaimaterials} explicitly couples a hypothesis-generation cognition module to an autonomous experimental-verification module driving physical synthesis robots, instantiating a complete generate–execute–feedback loop for materials. These systems make the survey's central thesis tangible: in a self-driving lab, the agent's policy \textit{is} the code, the lab \textit{is} the runtime, and the publication record \textit{is} the log.

\paragraph{Toward Agentic and Instruction-Following Science}

A final dimension of code-as-harness scientific agents is controllability: the ability to steer them with high-level scientific intent while preserving rigorous execution semantics. Benchmarks have rapidly emerged to measure this capability. MLAgentBench \citep{huang2024mlagentbenchevaluatinglanguageagents} evaluates language agents on 13 open-ended ML research tasks, requiring agents to read code, run experiments, and improve metrics. MLE-bench \citep{chan2025mlebenchevaluatingmachinelearning} scales this to 75 Kaggle ML-engineering competitions; the best-performing scaffold at release (OpenAI o1-preview with the Weco AIDE tree-search agent \citep{hu2025surveyscientificlargelanguage}) reaches Kaggle bronze-medal level on 16.9\% of competitions, and AIDE achieves roughly three times the medal rate of the next agent. ScienceAgentBench \citep{chen2025scienceagentbenchrigorousassessmentlanguage} compiles 102 tasks adapted from peer-reviewed publications across bioinformatics, computational chemistry, GIS, and cognitive neuroscience, \textit{unifying every target output as a self-contained Python program}, which is an explicit endorsement of code as the universal interface to data-driven science. DiscoveryBench \citep{majumder2024discoverybenchdatadrivendiscoverylarge} complements this with 264 multi-step hypothesis-search tasks across six domains, exposing failure modes of current agents (best system score $\sim$25\%). On the controllability side, instruction-following progress is visible in systems such as the AI co-scientist \citep{gottweis2025aicoscientist}, where scientists steer the multi-agent debate via natural-language research goals and constraints, in Biomni \citep{huang2025biomni}, whose graphical interface accepts natural-language queries and returns auditable code execution, and in the Virtual Lab \citep{swanson2025virtual}, where a human PI specifies high-level objectives and the AI PI dynamically configures a team of expertise-specific agents. AlphaEvolve \citep{novikov2025alphaevolvecodingagentscientific} and AlphaProof \citep{hubert2025olympiad} represent the goal-conditioned extreme: the agent is given only an objective function or a theorem statement, and the closed code-execution loop searches for any program that satisfies the verifier. Across these systems, instruction-following is realized by translating user goals into typed program specifications that the runtime can rigorously enforce.

Taken together, recent work on agents for scientific discovery exemplifies the survey's central shift: from static prediction toward interactive, stateful, and executable decision making. Hypotheses cease to be free-floating sentences and become parameterized programs; experiments cease to be lab notebooks and become version-controlled code; analyses cease to be one-off scripts and become reproducible artifacts that downstream agents can re-execute; and laboratories cease to be opaque physical sites and become production runtimes addressable through documented APIs. The result is a closed generate–execute–feedback loop in which a single substrate, code, carries scientific reasoning, scientific action, and the scientific environment itself, providing a unified foundation on which agents like the AI Scientist \citep{lu2024aiscientistfullyautomated,yamada2025aiscientistv2workshoplevelautomated}, AI co-scientist \citep{gottweis2025aicoscientist}, Virtual Lab \citep{swanson2025virtual}, Biomni \citep{huang2025biomni}, Coscientist \citep{boiko2023autonomous}, and AlphaEvolve \citep{novikov2025alphaevolvecodingagentscientific} can be compared, composed, and progressively improved. As benchmarks such as MLAgentBench \citep{huang2024mlagentbenchevaluatinglanguageagents}, MLE-bench \citep{chan2025mlebenchevaluatingmachinelearning}, ScienceAgentBench \citep{chen2025scienceagentbenchrigorousassessmentlanguage}, and DiscoveryBench \citep{majumder2024discoverybenchdatadrivendiscoverylarge} make precise, the open challenge is not whether code-as-harness agents can imitate isolated scientific tasks, but whether they can be trusted to drive the full loop autonomously, which is a challenge for which the program-world abstraction provides both the right ontology and the right experimental harness.

\subsubsection{Agent Personalization}

Personalization and recommender systems offer a distinctive setting for code-centric agentic systems. Unlike coding, GUI control, or scientific discovery, the environment here is not only a software system but also a human user whose intent, satisfaction, and long-term goals are only partially observed. As recommendation moves from static ranking toward interactive agents, the central challenge becomes how to maintain, update, and govern a user model through repeated interaction. Code is useful in this setting not simply because it executes recommendation policies, but because it provides an inspectable substrate for preference representation, feedback processing, constraint enforcement, and policy adaptation.

\paragraph{From Static Recommendation to Interactive Personalization}
Traditional recommender systems usually treat personalization as a prediction problem: given historical interactions, the system scores candidate items and returns a ranked list~\citep{he2020lightgcn,guo2017deepfm}. LLM-based recommenders broaden this view by enabling conversational preference elicitation, explanation, and multi-step refinement. Early prompting-based approaches query an LLM with user history and ask it to produce recommendations directly \citep{hou2024large, dai2023uncovering}. More agentic systems instead decompose recommendation into candidate retrieval, filtering, re-ranking, explanation, and feedback collection.
The emerging agentic recommendation~\citep{liu2025recoworld,wang2024recmind,huang2025recommender} instantiate this direction by using LLMs to coordinate recommendation sub-tasks through tool calls and structured intermediate states. Agent4Rec \citep{zhang2024generative} and iAgent~\cite{xu2025iagent} further simulates recommendation sessions with synthetic users, enabling offline evaluation of interactive policies. These systems mark a shift from recommendation as one-shot scoring to an adaptive process, where each interaction may revise the system's belief about the user.

\paragraph{Preference State as an Editable Artifact}
A key difference between personalization agents and other agentic systems is that the most important state is not fully observable. User preferences are latent, contextual, and often unstable. A user may click an item for convenience rather than genuine interest, skip an item because of timing rather than dislike, or change goals across sessions. Therefore, personalization agents need explicit preference states that can absorb noisy behavioral signals while remaining interpretable and correctable.
Code-centric representations provide a practical way to structure this state. Short-term interests can be stored as recent interaction logs, contextual summaries, or session-level preference vectors. Long-term preferences can be maintained as structured memory objects that record stable interests, constraints, and user-provided corrections. AMem \citep{xu2026mem} and related memory-based systems~\citep{wei2025evo,chhikara2025mem0} show how long-term user information can be maintained as editable documents or structured records. MemRec \citep{chen2026memrec} further studies how collaborative signals can support memory management for personalized recommendation. Compared with opaque embedding-only memory, structured preference memory is easier to inspect, revise, and reuse. A user can correct a stored preference in natural language, and the system can update the corresponding state before generating future recommendations.

\paragraph{Feedback as Policy Adaptation}
Personalization agents are driven by feedback, but the feedback is often sparse, delayed, and ambiguous. Clicks, dwell time, ratings, purchases, skips, and conversational corrections all provide partial evidence about user satisfaction. Production recommender systems already rely on code-defined feedback pipelines that log interactions, compute metrics, run A/B tests, and trigger model or policy updates. In an agentic setting, these pipelines become part of the personalization harness: they determine what signals are recorded, how they are interpreted, and when the agent should adapt.
User simulators~\citep{zhang2025llm,wang2025user,liu2025recoworld} provide an offline way to study such adaptation. They allow recommendation policies to be tested under controlled behavioral assumptions before real deployment. Recent LLM-based simulators extend this idea by generating richer synthetic user profiles and interaction traces. However, the central difficulty remains that simulated feedback may not match real user behavior, especially when recommendations themselves influence future preferences.

\paragraph{Controllable and Instruction-Following Personalization}
A major opportunity for agentic personalization is to move beyond optimizing implicit engagement signals toward following explicit user instructions. Users may want recommendations that satisfy constraints such as avoiding certain sources, limiting repeated categories, balancing exploration and familiarity, or prioritizing long-term goals over short-term engagement. These requirements are hard to express through a single learned score but can be represented as structured constraints, filters, or reward functions.
LLM-based conversational recommenders can elicit such preferences in natural language and translate them into policy specifications \citep{hou2024large}. Constraint-based recommendation further shows how fairness, diversity, and exposure requirements can be enforced at serving time rather than hidden inside model parameters \citep{lei2020conversational}. Explanation-based systems provide another path toward controllability: if a system explains why an item was recommended, the user can correct the rationale, and the corrected explanation can update the preference state. This makes personalization more interactive and auditable, since the user can shape not only outputs but also the logic behind future outputs.

\paragraph{Open Challenges for Personalization Harnesses}
Personalization raises several challenges that are sharper than in other domains. First, preference grounding remains unresolved. Unlike code assistants, which can rely on tests, or GUI agents, which can check interface states, personalization agents lack a reliable oracle for true user satisfaction. Proxy metrics such as clicks and engagement can be misleading or even harmful when optimized too aggressively.
Second, preference memory introduces privacy and governance risks. Long-term user models may contain sensitive behavioral patterns, so the harness must specify what is stored, where it is stored, how it is updated, and how users can inspect or delete it. Third, personalization is inherently multi-stakeholder. A platform may optimize engagement, a creator may seek exposure, and a user may value welfare or autonomy. Reducing these objectives to a single reward function can obscure conflicts of interest.



\subsection{Open Problems}

Code-as-harness systems shift the central challenge of agentic AI from isolated model generation to the reliability of the complete execution loop. Once agents act through tools, memory, code execution, shared state, and environment feedback, failures may arise from weak verifiers, stale context, unsafe tool access, inconsistent multi-agent state, insufficient multimodal grounding, or poorly governed self-improvement. These issues cannot be diagnosed by final task success alone. This section outlines the key open problems that emerge when the harness is treated as a first-class system component, with the goal of building agentic systems that are executable, inspectable, stateful, verifiable, and governed in long-horizon real-world environments.

\subsubsection{Harness-Level Evaluation and Oracle Adequacy}

Evaluation becomes difficult once an LLM is embedded in a code-agent harness. In this setting, performance is no longer determined by the base model alone, but also by the surrounding runtime: which repository files are retrieved, which tools are exposed, how many retries are allowed, whether the agent can execute tests, how failures are summarized, and what verifier decides success. However, most existing evaluations measure end-task success: whether a generated solution passes tests, solves an issue, or completes an interactive task. Such metrics conflate the capabilities of the base model, the quality of the harness, the reliability of tools, the informativeness of feedback, and the difficulty of the environment. This is especially visible in repository-level software engineering, where an agent may pass visible tests while exploiting weak or incomplete test suites; in GUI/OS tasks, where a scripted checker may miss unsafe or undesirable intermediate actions; and in scientific or embodied settings, where successful execution in a simulator may not imply that the result is scientifically valid or physically safe \citep{jimenez2024swebench,deng2025swe,miserendino2025swe,merrill2026terminal,jain2024livecodebench,chen2025scienceagentbenchrigorousassessmentlanguage}.

A key open problem is therefore to define \emph{harness-level metrics} that evaluate the operational substrate itself. These metrics should complement final task accuracy with measurements of execution reliability, feedback quality, context sustainability, safety, coordination, and reproducibility. Useful dimensions include: (i) \emph{trajectory efficiency}, such as number of tool calls, tokens, edits, executions, and wall-clock time; (ii) \emph{verification strength}, such as test coverage, oracle diversity, and rate of false acceptance; (iii) \emph{recovery ability}, such as whether the agent can diagnose and repair failures after invalid actions; (iv) \emph{state consistency}, such as whether memory, repository state, execution traces, and agent beliefs remain synchronized; (v) \emph{safety compliance}, such as whether permissions, sandboxes, and human-approval gates are respected; and (vi) \emph{replayability}, such as whether the full trajectory can be reconstructed and audited from logs and artifacts~\citep{anthropic2026agentevals}. A central bottleneck in this agenda is \emph{oracle adequacy}: whether the evaluator captures the intended task rather than only a narrow executable proxy. The open problem is not merely to build harder benchmarks, but to evaluate the code-agent harness as an executable runtime system.




\subsubsection{Semantic Verification Beyond Executable Feedback}

Oracle adequacy becomes especially challenging because execution feedback, while central to code-centric agents, can create a false sense of correctness: code can be run, traces can be inspected, tests can be checked, and failures can be fed back into revision. However, execution is only as reliable as the oracle attached to it. Unit tests may be incomplete, static analyzers may over-approximate, GUI checkers may miss unacceptable intermediate actions, scientific scripts may encode invalid assumptions, and robot simulators may hide physical risks. As a result, a harness can become overconfident precisely because it has executable feedback: the agent sees a green test, but the green test is not the full specification.

The central missing abstraction is a verification stack with explicit scope. Instead of treating pass/fail as a single terminal signal, future harnesses should compose multiple verification artifacts: unit tests, integration tests, property-based tests, fuzzers, static analyzers, type checkers, security scanners, runtime monitors, coverage reports, formal specifications, model-based critiques, and human review. Each artifact should declare what it verifies, what it cannot verify, and what confidence it provides. This is especially important for self-repair and self-evolving harnesses: if the verifier is weak, the agent will learn to optimize against the wrong signal. A useful direction is to make every accepted action carry an evidence bundle containing the checks run, the assumptions preserved, the untested regions, and the remaining risks. In this view, verification is not a final gate; it is an evolving, inspectable contract between the agent, the harness, and the environment.

Other promising directions include feedback calibration, independent verification, metamorphic testing, differential testing, property-based test generation, execution-trace summarization, and uncertainty-aware critics~\citep{ni2023lever,jung-etal-2025-code,tang2026execverify}. Reliable feedback should also be routed differently depending on its type: compiler errors may trigger local syntax repair, test failures may trigger behavioral diagnosis, coverage gaps may trigger test generation, and inconsistent reviewer comments may trigger arbitration. The broader goal is to build feedback loops that are not merely reactive, but epistemically aware: the harness should know when a signal is strong enough to act on, when it is weak, and when additional evidence is required.


\subsubsection{Self-Evolving Harnesses without Regression}

Most current harnesses are manually designed: developers choose the planning loop, memory format, tool set, permission rules, debugging procedure, and agent topology. However, as tasks become longer and more diverse, fixed harnesses may be suboptimal. A harness that works well for competitive programming may fail for repository repair; a harness tuned for GUI navigation may be inefficient for scientific workflows; and a multi-agent topology that succeeds on one task distribution may waste computation on another. This suggests that future systems should treat the harness itself as a programmable component that can adapt to new environments, rather than a fixed wrapper around the base model.

Automatic harness evolution is already underway. AutoHarness synthesizes code harnesses that constrain invalid actions~\citep{lou2026autoharness}, MetaHarness searches over harness code~\citep{lee2026metaharness}, Agentic Harness Engineering evolves harness components from observability signals~\citep{lin2026agentic}, and related methods optimize prompts, contexts, and workflows through reflection, search, or execution feedback \citep{agrawal2025gepa,Liu2025SEW,zhang2025agentic}. These systems point toward a broader paradigm in which an overarching optimization process analyzes runtime feedback, such as computational cost, decision paths, tool-use traces, memory pressure, and specific failure cases, and proposes modifications to the harness itself. Such modifications may reorganize communication among sub-agents, adjust memory allocation, revise retrieval or verification policies, or change how execution feedback is routed through the system. Therefore, ``automated harness evolution'' is not itself the open problem. The harder problem is whether a harness can improve itself without overfitting, weakening safety, increasing cost, hiding failures, or regressing on rare but important tasks.

The central insight is that a harness mutation should be treated like a code change to a safety-critical runtime. Every proposed edit should carry a change contract: which component is modified, which failure mode it targets, what improvement it predicts, which invariants it must preserve, which evaluation can falsify it, and how it can be rolled back. This is especially important because harness changes affect the future distribution of agent behavior. A new retrieval policy may improve benchmark accuracy while increasing hallucinated evidence; a new tool schema may reduce token cost while weakening permission boundaries; a new verifier may improve pass rate by accepting underspecified solutions. Future work should develop evidence-carrying harness evolution, held-out regression suites, safety invariants, canary deployment, rollback semantics, and causal evidence for why a harness edit helped. The goal is not a harness that changes often, but one that changes only when it can justify the change. A practical research agenda includes: defining mutation operators for harness components; building telemetry standards; evaluating evolved harnesses across diverse tasks; enforcing safety invariants during evolution; and separating improvements in the harness from improvements in the base model.


\subsubsection{Transactional Shared Program State and Semantic Conflict Resolution}

Scaling from single agents to multi-agent systems turns the codebase into a shared harness substrate. Planners, coders, testers, reviewers, security agents, and humans may all read and modify overlapping artifacts. Prior sections show that many systems still rely on sequential handoff, shared logs, or file-only state, while newer systems introduce blackboards, repository memories, execution feedback, and explicit belief-state synchronization \citep{Qian2023ChatDev,Hong2023MetaGPT,huang2023agentcoder,wang2025openhands,Guo2025SyncMind}. The open problem is that synchronization alone does not provide transactional semantics or assumption-level consistency: these mechanisms often synchronize artifacts but not assumptions. One agent may plan from an old repository snapshot, another may test a newer patch, a third may remember an obsolete invariant, and a human reviewer may introduce a new constraint that is not propagated to the rest of the system.

The missing abstraction is transactional shared program state. Agents should not merely append messages to a common log; each action should declare its read set, write set, assumptions, version dependencies, verifier obligations, and conflict policy. Conflicts should be detected not only at the level of file diffs, but also at the level of plans, tests, retrieved evidence, permissions, memory entries, and latent user requirements. Future harnesses need conflict-resolution mechanisms that are semantic rather than purely textual, including semantic merge, rollback, dependency-aware locking, belief-state reconciliation, conflict explanation, and re-verification after merge. Classical version control, databases, CRDTs, and build systems provide useful analogies, but agentic systems add conflicts that conventional tools do not see: incompatible plans, stale memories, duplicated subtasks, inconsistent tool authority, and divergent interpretations of the user's goal. A key research challenge is to determine when a conflict can be resolved automatically and when it requires external judgment. Such mechanisms also require metrics beyond merge correctness, including merge success, semantic regression rate, rollback frequency, conflict recurrence, and the cost of human intervention.


\subsubsection{Human-in-the-Loop Safety and Accountability as Harness State}

As code-as-agent-harness systems are used in increasingly consequential settings, safety cannot be delegated to the base model or encoded only as a natural-language instruction. In critical domains such as software deployment, cybersecurity, finance, healthcare, scientific experimentation, enterprise automation, and embodied control, agent actions may affect production systems, private data, external users, physical devices, or institutional compliance. A harness therefore needs to function not only as a context manager or tool executor, but also as a safety governor between model intent and real-world consequence. It should classify proposed actions by risk, enforce permission tiers, deny actions that violate hard constraints, and require human approval for irreversible or externally consequential transitions. For example, when an agent requests credentials, modifies security-critical code, accesses user data, deploys a service, issues financial or medical recommendations, or controls physical equipment, the harness should be able to override the base model and suspend autonomy until a human decision is made \citep{Nunez2024AutoSafeCoder,vijayvargiya2025openagentsafety,guan2025normcode}.

Future harnesses need explicit governance mechanisms that mediate between model intent and environmental action. A useful design pattern is a multi-tier permission model. At the lowest tier, agents may read files, inspect logs, and run static analysis. At higher tiers, they may edit local files, execute sandboxed code, access the network, call external APIs, modify shared repositories, or affect production systems. Each tier should specify its allowed actions, constraints, audit logs, rollback mechanisms, and human-in-the-loop gates for high-risk operations. Such governance must also be context-sensitive. The same command may be safe in a disposable sandbox but unsafe in a production repository, and the same network request may be benign during documentation retrieval but risky when it transmits local state. Therefore, permissions should depend not only on tool identity, but also on arguments, environment state, data sensitivity, and expected side effects. Open problems include policy specification, side-effect prediction, sandbox escape prevention, secret handling, secure tool schemas, reversible execution, and measuring the tradeoff between autonomy and safety.


This safety role also changes how human feedback should be represented. Human-in-the-loop control should not appear only as an occasional prompt interruption; it should become durable harness state. Each approval, rejection, policy exception, or reviewer correction should update the harness's permission rules, escalation policy, verification criteria, and future memory retrieval. Likewise, high-stakes approvals should be auditable state transitions: what action was proposed, what evidence was shown, what risks were surfaced, who approved or rejected it, and what responsibility boundary changed afterward. The open problem is to design harnesses that can decide when autonomy is appropriate and when human judgment is mandatory. In this view, reliable code-as-agent-harness systems require not only executable code and verifiable feedback, but also executable accountability: a safety layer that filters, vetoes, escalates, and records agent actions before they reach the real world.

\subsubsection{Multimodal Code-Harness Systems}

Most code-agent harnesses are still designed around textual state: prompts, files, logs, tool outputs, tests, and execution traces. However, many emerging agentic systems operate in environments where the critical state is multimodal. GUI agents observe screenshots, accessibility trees, and rendered interface states; embodied agents rely on egocentric images, depth, force, tactile signals, object poses, and simulator or robot states; scientific agents inspect plots, microscope images, molecular structures, and experimental readouts. In these settings, the harness can no longer treat perception as a passive input to the model. It must manage multimodal observations as persistent, queryable, and verifiable state.

A central challenge is multimodal context compression. Visual observations are large, redundant, and often only partially relevant to the task. A GUI screenshot may contain hundreds of elements, while only one button matters; an embodied trajectory may contain thousands of frames, while only a few reveal task-critical object relations, contact events, or failure causes. Future harnesses need compression mechanisms that preserve task-relevant visual evidence rather than merely reduce token cost. This suggests a multi-level memory design: raw images or frames are stored as immutable evidence; object-, region-, element-, and pose-level annotations provide structured intermediate state; and compact textual or symbolic summaries expose only the information needed for skill retrieval and planning. The open problem is to decide what multimodal information should be retained, abstracted, forgotten, or promoted into long-term memory, especially when later failures reveal that an earlier visual or physical detail was important.

Visual grounding introduces a second challenge: aligning observations with actions. In text-centric harnesses, an action can often be checked against a file, command, or test result. In visual environments, the agent must map language goals to image regions, interface elements, objects, coordinates, poses, and executable actions. A GUI agent must know that a planned click corresponds to the correct rendered button; an embodied agent must know that a grasp command targets the intended object under the current camera view and physical configuration. This requires harness-level grounding contracts that connect perception, action, and verification. Each action should carry not only a natural-language rationale, but also a grounded reference to the evidence it depends on, such as a bounding box, object identifier, UI element, frame index, region feature, object position, or orientation. After execution, the harness should verify whether the intended grounded state changed as expected, rather than relying only on the model's self-report.

Reliable feedback is also harder in multimodal settings. A textual error message or unit-test failure provides an explicit signal, but visual and physical feedback is often implicit, delayed, or ambiguous. A button may look clicked without triggering the right state transition; a robot may appear to hold an object while the grasp is unstable; a chart may seem to support a conclusion while its axis scale changes the interpretation. Future harnesses therefore need multimodal verification stacks that combine visual state checks, object tracking, OCR or UI-tree inspection, simulator state, physical sensors, tactile feedback, and task-specific validators. More importantly, each feedback signal should expose its scope and uncertainty. For example, a bounding-box detector verifies localization but not task completion; a simulator state verifies object position but not physical robustness; an OCR result verifies visible text but not semantic correctness. This also calls for tighter integration between world modeling and action modeling: the harness should predict how the visual or physical world is expected to change after an action, compare that prediction with the observed outcome, and use the mismatch to diagnose failures. In embodied and robotic settings, such prediction-error signals are especially important for recovery, since failures may arise from occlusion, slippage, collision, unreachable poses, or violated preconditions rather than from an explicit error message. Treating multimodal feedback as calibrated evidence, rather than as a binary success signal, is essential for safe long-horizon autonomy.

Multimodal memory should also support skill evolution. In visual-centric domains such as GUI control and embodied manipulation, reusable skills cannot be represented only as text or code snippets. A useful skill often couples a multimodal precondition, an executable action pattern, and an expected postcondition: what the agent should see or sense before acting, what program, UI command, or motor primitive it should execute, and what visual, physical, or state change should follow. For example, a GUI skill may encode how to locate a settings menu from a screenshot, click the correct region, and verify that a new panel appears. An embodied skill may encode how to identify a graspable object, choose an approach pose, execute a primitive controller, and confirm through vision, force, or tactile feedback that the object has moved into the gripper. Such skills should evolve from successful trajectories, failed attempts, and human corrections, while retaining their grounding evidence. The harness must therefore decide when a visual-action pattern is reusable, how abstractly it should be stored, and how to adapt it across layouts, viewpoints, embodiments, sensors, or tasks.

\subsubsection{Toward a Science of Harness Engineering}

Taken together, these open problems suggest that code-as-harness research is moving toward a broader science of harness engineering. The central object of study is no longer only the model or the generated program, but the complete closed-loop system: context, memory, tools, execution, feedback, safety, coordination, and evaluation. Progress will require benchmarks that expose long-horizon failures, telemetry that makes trajectories auditable, metrics that isolate harness components, and design principles that allow agents to operate safely in persistent program worlds.

The most important future systems will likely be those that combine four properties. First, they will be \emph{executable}, grounding decisions in code, tools, tests, and environments. Second, they will be \emph{inspectable}, exposing plans, state, provenance, and failure causes. Third, they will be \emph{stateful}, preserving task-relevant information across long trajectories and multiple agents. Fourth, they will be \emph{governed}, ensuring that autonomy is constrained by permissions, verification, and accountability. These properties define the next frontier for reliable, long-horizon agentic AI.
